Los resultados permiten contrastar las dos hipótesis formuladas en el estudio. En relación con la primera, que planteaba una precisión cercana al $70\%$, los valores obtenidos muestran un comportamiento inferior al esperado cuando la capacidad predictiva se evalúa mediante el coeficiente de determinación ($R^2$). El mejor resultado se observó en la competencia de Inglés, con un valor de $0,61$, mientras que las demás competencias registraron valores comprendidos entre $0,39$ y $0,51$. Respecto a la segunda hipótesis, los resultados tampoco evidencian una mejora promedio del $10\%$ entre las pruebas Saber 11 y Saber Pro. Por el contrario, la medición del valor agregado institucional muestra resultados negativos en la mayoría de las competencias y años analizados. Este hallazgo indica que el desempeño alcanzado por los estudiantes fue, en promedio, menor al que el modelo estimaba a partir de sus condiciones iniciales.
Conviene precisar por qué un coeficiente de determinación moderado no constituye, en sí mismo, una falla del sistema, sino una consecuencia esperable del problema. Entre la presentación de Saber 11 y la de Saber Pro median alrededor de cinco años, durante los cuales ocurre precisamente aquello que el estudio busca medir: el efecto de la formación universitaria. La porción del desempeño en Saber Pro que no puede explicarse desde Saber 11 es, en buena parte, el valor agregado de la institución; si un modelo lograra explicar la totalidad de la varianza, ello implicaría que la universidad no aporta nada y que el resultado final ya estaba determinado por las condiciones de entrada. En este marco, una capacidad explicativa moderada es coherente con la naturaleza del fenómeno y se alinea con la literatura de minería de datos educativos, que reporta capacidades predictivas igualmente acotadas en problemas análogos.
Por esta misma línea de razonamiento, el periodo temporal transcurrido entre la presentación de ambas pruebas constituye un intervalo de incertidumbre, donde el desarrollo académico de una persona no está condicionado únicamente por su proceso académico en la institución universitaria. La mejora o desmejora de las capacidades y competencias de un alumno están, en efecto, afectadas por la Institución de Educación Superior a la cual pertenece, sin embargo, de igual forma influyen una cantidad de parámetros ajenos a la universidad que no es posible prever en su totalidad y por tanto este estudio no contempla. Un alumno A podría mantener hábitos de lectura como actividad de ocio mientras que un alumno B no, este alumno B podría frecuentar actividades extracurriculares de estudio por el aspecto social (que fomentan la mejora de sus capacidades) mientras que un alumno C no; es este aspecto personal, e imprevisible en gran parte, que no se ve representado en la información recopilada en los datasets de estudio y por tanto genera una variabilidad en el modelo que implica una gran incertidumbre en el sistema al momento de intentar predecir puntajes de competencias.
El segundo hallazgo relevante es que la incorporación de variables socioeconómicas no incrementó la capacidad predictiva, e incluso la redujo en dos competencias. Una explicación plausible es que el efecto del contexto socioeconómico ya se encuentra incorporado de manera implícita en los propios puntajes de Saber 11, de modo que añadir esas variables introduce más ruido que señal. Este comportamiento ilustra empíricamente el principio de la maldición de la dimensionalidad descrito por \cite{verleysen2005curse}: ampliar el vector de entrada de cuatro a veinte variables no se tradujo en una mejora, lo que sugiere que la dimensionalidad adicional no estuvo acompañada de información discriminante nueva.
El caso del Inglés merece atención particular, pues es la única competencia con valor agregado positivo en los primeros años y la única en la que un modelo lineal iguala a la red neuronal. Ambos hechos son consistentes entre sí y con la evidencia de que el dominio del idioma está fuertemente condicionado por el nivel socioeconómico: se trata de la competencia más predecible desde las condiciones de entrada y, a la vez, aquella en la que la institución parece incidir de forma más visible. Este resultado respalda la decisión metodológica de calcular el puntaje global de Saber Pro mediante un promedio ponderado que otorgue un menor peso a la competencia de Inglés, siguiendo un criterio similar al que utiliza el ICFES en la prueba Saber 11. 
De esta manera, se busca que la ponderación refleje de forma más equilibrada la contribución de cada competencia al resultado global. Por otra parte, los resultados relacionados con la disminución del valor agregado observada a partir de 2019 deben analizarse teniendo en cuenta el contexto en el que ocurrieron. Este periodo coincide con la pandemia de COVID-19, una situación que afectó de manera significativa las condiciones de enseñanza, aprendizaje y evaluación. Por esta razón, los cambios observados no pueden asociarse exclusivamente al desempeño institucional, ya que también estuvieron influenciados por factores externos que impactaron el proceso educativo. Asimismo, conviene recordar que la partición temporal de los datos (entrenar con cohortes pasadas y evaluar con una posterior) ofrece una estimación honesta, y por tanto más conservadora, del desempeño que la que arrojaría una partición aleatoria.